\documentclass[11pt]{article}
\usepackage[margin=1in]{geometry}
\usepackage{amsmath,amssymb}
\usepackage{xcolor}
\usepackage{enumitem}
\usepackage{array}
\usepackage{longtable}
\usepackage{booktabs}
\usepackage{hyperref}
\usepackage[normalem]{ulem} % \uline: an underline that breaks across lines
\hypersetup{colorlinks=true, linkcolor=blue, urlcolor=blue}

% Reviewer comment vs. our response formatting.
\newcommand{\comment}[1]{\medskip\noindent\textbf{Comment.}\ \textit{#1}\par}
\newcommand{\response}[1]{\smallskip\noindent{\color{blue!60!black}\textbf{Response.}}\ #1\par\medskip}
\newcommand{\loc}[1]{\textcolor{blue!50!black}{[#1]}}

\newenvironment{point}[1]{\medskip\noindent\textbf{\uline{#1}}\par}{\par}

\title{\vspace{-2em}Response to Reviewers\\[0.3em]
\large Manuscript: \emph{A Dual-Track Adaptive Projection Algorithm for a Nonlinear
Split Feasibility Problem with Mixed Convex--Pseudoconvex Constraints}\\[0.2em]
\normalsize Journal: \emph{Optimization Methods and Software}}
\author{Nguyen Van The, Tran Ngoc Thang, and Nguyen Hong Son}
\date{}

\begin{document}
\maketitle
\vspace{-2em}

\noindent We thank the Editor and both Reviewers for their careful and constructive
reading. The comments have materially strengthened the paper, and in several places led
us to remove overstated claims and to add new, more honest evidence. Below we respond to
each point individually and indicate the corresponding change in the revised manuscript.
\textbf{All added or modified text in the revised manuscript is typeset in
\textcolor{red}{colored text}} so the changes are easy to locate. Section, equation,
theorem, and table numbers below refer to the revised manuscript. In summary, the
revision:
\begin{itemize}[nosep]
  \item merges the former Introduction and Related work into a single, shorter
        Introduction and cuts the paper from $39$ to $25$ pages, rewriting it throughout
        in shorter sentences;
  \item qualifies the ``constant-free'' claim (new Remark~2) and reframes the proportional
        weight as a heuristic residual-share allocation, with no claim of necessity or
        optimality (new Remark~1);
  \item makes the whole-sequence convergence argument explicit (rewritten Stage~3 of
        Theorem~3.12), clarifying that it needs only one feasible cluster point and
        Fej\'er monotonicity, not a unique limit;
  \item adds a solver diagnostic (new Sec.~4.4) reporting per-trajectory constraint
        activation, a tuned-modulus subgradient comparison with its tuning cost, and
        controlled instances with a genuinely nonconvex pseudoconvex constraint;
  \item proves that the feature-selection ratio is in fact convex (new Proposition~4.1),
        states plainly that the application is a convex special case, and removes all
        claims of uniform superiority;
  \item unifies the reference format, adds the requested self-adaptive and decomposition
        works for the SFP, moves the appendix results into the main text, and resolves
        every notational and presentation issue raised.
\end{itemize}

% =====================================================================
\section*{Response to the Editor}

\begin{point}{The Associate Editor recorded no specific comments; the decision was major
revision.}
\response{We have revised the manuscript in full accordance with both Reviewers' reports,
as detailed below, and have highlighted all changes in colored text. The two most
substantial structural changes --- merging and shortening the front matter (Reviewer~2)
and adding the solver diagnostic with constraint-activation and tuned-baseline evidence
(Reviewer~1) --- are summarized in the list above.}
\end{point}

% =====================================================================
\section*{Response to Reviewer \#1}

\noindent We thank the Reviewer for the accurate summary of the contribution and for five
comments that led directly to a more careful statement of the theory and a more complete
experimental section.

\begin{point}{R1.1 The soft weight $\mu_k=J_k^{\mathrm{cvx}}/(J_k^{\mathrm{cvx}}+J_k^{\mathrm{ps}})$
is central to the algorithm's novelty, yet the convergence proof only requires $\mu_k$ to
satisfy certain summability conditions -- any sequence in $[0,1]$ would suffice. The chosen
form is therefore not structurally necessary. The authors should either justify this
specific weighting or reframe the contribution accordingly.}
\response{We have adopted the Reviewer's second suggestion and reframed the contribution.
We agree that the proportional form is not structurally necessary: the one-step
quasi-Fej\'er inequality (Lemma~3.8) holds for every $\mu_k\in[0,1]$, and the convergence
proof needs only that the weight keep each track active --- with the constant weight
$\mu_k\equiv\tfrac12$, for instance, the proofs of Lemma~3.10 and of Stage~2 go through with
only minor changes. The experiments point the same way: a balanced variant, which assigns
equal weights when both tracks are violated, succeeds in $26/36$ runs on the feature
instances and $24/36$ on the nonconvex family, the same counts as DTAP \loc{Sec.~4.4,
Table~7}. We therefore no longer present the specific weight as a theoretical necessity or
an optimum. The Introduction now states that ``the proof uses this weight only to keep a
persistently violated track active''. New Remark~1 \loc{Remark~1, Sec.~3.4} states this explicitly: it records where
the proportional form enters the present proof and notes that the one-step inequality,
although valid for every $\mu_k\in[0,1]$, does not by itself imply convergence for every
weight sequence. The degenerate regimes of Corollaries~3.13--3.14 correspond to
$\mu_k\equiv1$ and $\mu_k\equiv0$. The dual-track
construction --- a residual-scaled step where a quantitative gradient inequality is
available and a safeguarded step where only sign information is --- remains the
contribution; the particular weight that blends them does not.}
\end{point}

\begin{point}{R1.2 The ``constant-free'' claim is overstated. The algorithm does not
explicitly require Lipschitz constants, but Assumption~1(A2) and Lemma~4.9 ensure the
iterates lie in a compact set where a Lipschitz constant necessarily exists. The claim
should be qualified.}
\response{We agree. Such an implicit dependence is shared by any method with bounded
iterates and continuous gradients, which is why we now frame the point as being about
algorithm \emph{inputs} rather than about the existence of a constant. The precise statement is now that no
operator norm, Lipschitz constant, or H\"older modulus is an \emph{input} of the
iteration, while the \emph{analysis} does use finite bounds on function values and
gradients over the compact set that contains the (provably bounded) iterates. New
Remark~2 \loc{Remark~2, Sec.~3.4} makes this explicit: ``Bounded gradients imply
Lipschitz continuity of the functions on that convex compact set, but these bounds are
not algorithm inputs,'' and notes that boundedness is established first (Lemma~3.9, the
former Lemma~4.9) before any such bound is invoked. The abstract now says that the
step sizes ``require no prescribed operator norm, Lipschitz constant, or H\"older modulus''
and that ``the proof uses local bounds on the bounded iterate region'', and the Conclusion
that ``no regularity modulus is an input, although the analysis uses bounds on a compact
set containing the iterates''. We no longer use the unqualified phrase ``constant-free.''}
\end{point}

\begin{point}{R1.3 Theorem~4.12 claims convergence of the whole sequence to a feasible
point. However, Stage~2 of the proof only establishes that at least one limit point is
feasible. Since the problem is nonconvex, multiple limit points may exist. The authors
should either strengthen the argument or revise the theorem to subsequential
convergence.}
\response{We have kept the whole-sequence statement, because it is correct, and rewritten
Stage~3 so that the step the Reviewer flagged is now explicit. The argument does not
require a unique limit point. For \emph{every} fixed $x^\star\in\Omega$, Lemma~3.9(i)
shows that $\{\lVert x^k-x^\star\rVert^2\}$ converges (Fej\'er monotonicity with respect
to the whole feasible set, using only $R^{x^\star}_k\ge0$ from Lemma~3.5). Stage~2
produces \emph{one} feasible cluster point $\bar x\in C\cap\Omega$. Applying Lemma~3.9(i)
with the admissible choice $x^\star=\bar x$, the limit $\lim_k\lVert x^k-\bar x\rVert^2$
exists; since a subsequence converges to $\bar x$, that limit is $0$, so the \emph{whole}
sequence converges to $\bar x$. Multiple feasible cluster points therefore cannot coexist.
The rewritten Stage~3 states this in full and closes with: ``The argument does not
require $\Omega$ to be a singleton: one feasible cluster point and the convergence of its
distance suffice'' \loc{proof of Theorem~3.12,
Sec.~3.8}. We have also aligned the Introduction (contribution~C2) with this wording.}
\end{point}

\begin{point}{R1.4 The thresholds $b_i$ are set as percentiles of sampled $F_i(w)$ to
guarantee $\Omega\neq\emptyset$. In real applications, thresholds are given by the problem,
not adjustable. More critically, if $b_3$ is so large that $F_3(w)\ll b_3$ throughout, the
pseudoconvex constraint is never binding. The authors should report constraint activation
frequencies and confirm the mixed setting is genuinely engaged.}
\response{We now address both halves of this comment directly, in the new subsection
\emph{Constraint activity and a tuned subgradient comparison} \loc{Sec.~4.4}. On
thresholds: we clarify that externally specified thresholds remain fixed inputs, and that
the percentile calibration is used only to define the benchmark instance, with the same
thresholds shared by all solvers; consistency ($\Omega\neq\emptyset$) is treated as an
assumption, and the exact levels are recorded in the supplementary metadata
\loc{Sec.~4.1.3 and Sec.~4.4}. On activation: we report the frequency of
$F_3(w^k)-b_3>10^{-6}$ along every trajectory. The third constraint is in fact
\emph{inactive} on seven of the eight datasets and active at only $1$ of $48$ evaluated
points on Abalone (Table~6). We state the consequence plainly, agreeing with the
Reviewer: on these instances the pseudoconvex track is essentially never binding, so the
dual-track mechanism is not needed there --- these feature-selection instances exercise a
convex subclass. Indeed, new Proposition~4.1 proves that the
fractional criterion $F_3=w^{\top}Q_Rw/\rho^{\top}w$ is \emph{convex} on a neighbourhood
of the simplex when $Q_R\succeq0$ and $\min_k\rho_k>0$, so the feature model is an
all-convex normal-box NSFP. To engage the mixed setting genuinely, we added controlled
instances with a provably nonconvex pseudoconvex constraint,
$F_3(x)=\kappa\log(1+\lVert x-e_1\rVert^2)$ \loc{Eq.~(47), Sec.~4.4}; there the
pseudoconvex track is active in all $36$ runs (pooled activity $12\,835/12\,877$ evaluated
points), so the pseudoconvex track is exercised directly rather than vacuously.}
\end{point}

\begin{point}{R1.5 DTAP and SG-const/SG-dyn are statistically indistinguishable on three
of four metrics (Table~5). SG-dyn's failure on $3$ datasets is attributed to poor
H\"older modulus estimation --- but with a tuned modulus it might perform comparably. An
additional experiment with manually tuned moduli would strengthen the comparison.}
\response{We added exactly this experiment \loc{Sec.~4.4, Tables~6--8}. For SG-dyn we now
report, alongside the sampled-modulus variant, a \emph{tuned} variant whose modulus is
selected from a manual grid $c\in\{10^{-3},\dots,10^{1}\}$ on held-out starts, and we
account for its tuning cost separately ($179\,744$ optimization evaluations and
$10.80$\,s on the feature instances; comparable figures for the synthetic family). Tuning
raises SG-dyn's standard-start success from $5/9$ to $9/9$ and, as the Reviewer
anticipated, a tuned subgradient method can match or outperform DTAP --- e.g.\ on Abalone,
and on the nonconvex family tuned SG-const and tuned SG-dyn reach the tolerance in all
$36$ runs versus $24/36$ for DTAP (Table~7). We have accordingly revised every comparative
claim: the text now states that no pairwise comparison is significant after Holm
adjustment, that the classification results reflect a size--prediction trade-off rather
than superiority, and, in the Conclusion, that ``tuned subgradient methods can match or outperform
DTAP''. The
reported advantage of DTAP is narrowed to the fact that no modulus has to be estimated or
tuned as an input of the iteration.}
\end{point}

% =====================================================================
\section*{Response to Reviewer \#2}

\noindent We thank the Reviewer for the detailed, page-by-page reading. We have followed
all comments; the paper is now substantially shorter and, we hope, considerably clearer.

\subsection*{Major issues}

\begin{point}{M1. Sections~1 (Introduction) and~2 (Related work) should be reorganized and
merged into Section~1, limited to about $4$ pages; the paper should be shortened (about
$25$ pages) by removing redundant content and long sentences.}
\response{Done. The Related work section (and its Table~1) has been removed and its
essential content condensed into the Introduction \loc{Sec.~1}, which now recalls the
$CQ$ method, self-adaptive and decomposition strategies, and the closest nonlinear
antecedents (Li et al., Yen et al., Hu et al.) in a few short paragraphs and states three
contributions instead of five. The Preliminaries were reduced to the facts actually used
and no longer carry subsections; the feature-selection framework (former Section~5) is now
a subsection of the experiments \loc{Sec.~4.1}; the sparse variant DTAP-S, the runtime
comparison, the robustness/ablation section, and the ``Summary of findings'' were removed;
the Conclusion is a single paragraph. The paper is now $25$ pages, down from $39$.}
\end{point}

\begin{point}{M2. The writing should be substantially improved: many sentences are
ambiguous or too long. Describe the problem, algorithms, etc.\ in short sentences; polish
each sentence; remove ``nonsense'' descriptions.}
\response{We rewrote the manuscript with this request in mind. Long sentences were split,
ambiguous phrases were replaced or deleted, and descriptive passages that did not support
a result were removed. The abstract, Introduction, Preliminaries, the opening of Section~3,
and the entire experimental section were rewritten. The parenthetical names of lemmas,
propositions, remarks, and corollaries were removed to shorten the text.}
\end{point}

\begin{point}{M3. The reference format is inconsistent, and important SFP references are
omitted; e.g.\ the authors of [23] and [35] also have self-adaptive or decomposition work
on the SFP.}
\response{We have unified the reference format to one style (initials, italic journal
names, volume, issue, pages, year, and DOI where available), and we now cite the
self-adaptive and decomposition works of the authors of [23] (Z.~Li, D.~Han, W.~Zhang) and
[35] (Q.~Yang) on the (multiple-sets) SFP: Zhang--Han--Li (\emph{Inverse Problems}, 2009),
Zhang--Han--Yuan (\emph{Comput.\ Optim.\ Appl.}, 2012), Zhao--Yang (\emph{Inverse
Problems}, 2011), and Yang--Yang--Zhang (\emph{J.\ Optim.\ Theory Appl.}, 2014). They are
cited in the first paragraph of the Introduction, where self-adaptive steps and
decomposition strategies for the SFP are recalled. References that were needed only for the
removed material were deleted.}
\end{point}

\begin{point}{M4. Why are some numerical results presented in the Appendix?}
\response{We agree, and the appendix has been removed. The per-dataset accuracy and AUC
values are now in the main text \loc{Table~3, Sec.~4.3}, and the two baseline solvers are
described in the main text \loc{Sec.~4.2.1}. The threshold and runtime tables were removed;
the thresholds of every instance are recorded in the supplementary reproducibility
package.}
\end{point}

\subsection*{Minor issues}

\noindent{\color{blue!60!black}\textbf{Response.}}\ We agree with all twenty minor
comments and have adopted every one of them without exception. For conciseness we list
them, with the corresponding change and its location, in the table below; in each case the
change is marked in colored text in the revised manuscript.

\medskip
{\footnotesize
\setlength{\tabcolsep}{5pt}
\renewcommand{\arraystretch}{1.25}
\begin{longtable}{@{}>{\raggedright\arraybackslash}p{0.42\textwidth}
                    >{\raggedright\arraybackslash}p{0.52\textwidth}@{}}
\toprule
\textbf{Reviewer's minor comment} & \textbf{Change made (location in revised manuscript)}\\
\midrule
\endfirsthead
\toprule
\textbf{Reviewer's minor comment} & \textbf{Change made (location in revised manuscript)}\\
\midrule
\endhead
\bottomrule
\endfoot
m1 (p.~2, l.~18) --- ``reference point'' is ambiguous.
  & Phrase removed; the L\'opez et al.\ rule is now cited among the self-adaptive methods in Sec.~1.\\
m2 (p.~2, l.~27) --- ``differing'' vs.\ ``different densities''.
  & ``Different''; the radiotherapy sentence was removed (Sec.~1).\\
m3 (p.~2, l.~35) --- polish ``The NSFP is not itself new''.
  & Sentence removed; Li et al., Yen et al.\ and Hu et al.\ are now compared with DTAP directly (Sec.~1).\\
m4 (p.~2, l.~39) --- ``active index set'' vs.\ ``active set''.
  & ``Active set''; sentence removed; we use ``approximate active set'' (Sec.~3.2).\\
m5 (p.~2, ll.~47--49) --- ``We ask instead \dots'' is unclear.
  & Removed; the class relation is restated in short sentences (Sec.~1).\\
m6 (p.~3, ll.~48--53) --- sentences too long.
  & Paragraph removed; reduced to one sentence (Sec.~1).\\
m7 (p.~3) --- five contributions too detailed.
  & Replaced by three short items C1--C3 (Sec.~1).\\
m8 (p.~5) --- ``Rule (3) has since been extended'' grammar.
  & Sentence removed with the Related work section.\\
m9 (p.~5, l.~33) --- ``drives their proof is for free''; ``nonconvexity sits in the sets''.
  & Unclear passage removed with the Related work section.\\
m10 (p.~5, l.~35) --- ``plain gradient inequality''.
  & Removed (it meant the convex gradient inequality $f(y)\ge f(x)+\langle\nabla f(x),y-x\rangle$).\\
m11 (p.~5, l.~44) --- ``the closest predecessor''.
  & Removed; the three antecedents are compared with DTAP directly (Sec.~1).\\
m12 (p.~5, l.~48) --- ``step'' vs.\ ``step size''.
  & ``Step size''; sentence removed with the Related work section.\\
m13 (p.~7, l.~31) --- ``dtap'' should be ``DTAP''.
  & Corrected throughout; the acronym is now always ``DTAP''.\\
m14 (p.~7, l.~39) --- ``Combettes 10'' should be ``Combettes~[10]''.
  & Corrected; the surviving citation reads ``in the sense of Combettes~[5]'' (Sec.~3.7).\\
m15 (p.~9) --- subsections 3.1--3.2 unnecessary; 3.3 should move to the experiments.
  & Preliminaries de-subsectioned (Sec.~2); the feature-selection model is now Sec.~4.1.\\
m16 (p.~12) --- $H:=\mathbb{R}^n$ repeated.
  & Repetition removed from the problem statement (Sec.~3.1).\\
m17 (p.~13) --- meaning of ``standing assumption''.
  & Subsection retitled ``Assumptions'' (Sec.~3.3); the name was dropped.\\
m18 (p.~18, l.~40) --- AMS style for $\sum_{k=0}^{\infty}\lVert x^{k+1}-x^k\rVert^2<\infty$.
  & Series displayed with limits above/below (Prop.~3.11); other displays checked.\\
m19 (p.~20, l.~36) --- polish ``so no Young splitting is needed''.
  & Phrase removed from the proof of Corollary~3.13.\\
m20 (p.~20, l.~35) --- expression $K\cap\bigcap_{i\in I^{\mathrm{cvx}}}$ is not good.
  & Rewritten as $\Omega=\{x\in K:F_i(x)\le b_i\ \forall i\in I^{\mathrm{cvx}}\}$ (Cor.~3.13), and as $\bigcap_{i\in I^{\mathrm{cvx}}}\{x\in K:F_i(x)\le b_i\}$ in the proof of Thm.~3.12.\\
\end{longtable}}

\begin{point}{General. Please consider the above comments and make the paper concise.}
\response{We have followed every comment and shortened the paper from $39$ to $25$ pages
(see M1 and M2).}
\end{point}

% =====================================================================
\section*{Additional changes made by the authors}

\noindent While revising, we also made the following changes, not requested by the
Reviewers, which we list for completeness and transparency.
\begin{enumerate}[leftmargin=1.6em,itemsep=3pt]
\item \textbf{Convexity of the fractional criterion.} The submitted version called
  $F_3(w)=w^{\top}Q_Rw/\rho^{\top}w$ pseudoconvex. We now prove it is \emph{convex} on an
  open neighbourhood of the simplex when $Q_R\succeq0$ and $\min_k\rho_k>0$
  (Proposition~4.1), so the feature-selection instance is an all-convex NSFP. The
  assignment $F_3$ to the pseudoconvex track used in the experiments remains admissible
  (convex $\Rightarrow$ pseudoconvex), and the all-convex assignment is reported as a
  comparator.
\item \textbf{Domain of the component functions.} The functions are now defined on an open
  neighbourhood $U$ of $K$, and (A2) requires continuous differentiability on $U$; this
  covers fractional functions with denominators positive on $K$. A short paragraph after
  Assumption~1 explains this and why $\Omega$ is convex.
\item \textbf{Scope of the analysis.} Besides Remarks~1 and~2 (see R1.1, R1.2), Remark~5
  now states precisely that the $O(N^{-1/3})$ bound concerns only the \emph{best}
  convex-track violation --- not a last-iterate bound, not a rate for the full residual or
  the distance to feasibility.
\item \textbf{Experimental protocol.} Section~4.2.2 now states the protocol exactly:
  min--max scaling is fitted once on the training pool, the voted masks are computed once
  and reused, the SVM probability calibration uses a fixed seed, and the evaluation is
  exploratory rather than nested. The Nogueira stability index (with the correct
  $N_f/(N_f-1)$ variance correction) and the statistical procedure (exact signed-rank
  test, Holm adjustment, Friedman test) are given explicitly.
\item \textbf{Recomputed values.} The stability values, the SVM AUC values, and the
  $p$-values were recomputed and now include raw and Holm-adjusted $p$-values; columns for
  SG-const and SG-dyn were added. Accuracies and all $5$-NN results are unchanged.
\item \textbf{Revised claims.} Statements of superiority (accuracy, speed, robustness)
  were removed or qualified: no pairwise comparison is significant after Holm adjustment,
  the results show a size--prediction trade-off, and tuned subgradient methods can match or
  outperform DTAP.
\item \textbf{Presentation of the algorithm.} The stopping test of Algorithm~1 is now
  stated at the start of Step~1 instead of in a separate Step~3; the iterates and the
  returned point are unchanged.
\item \textbf{Denominator safeguard.} Section~4.1.3 now reports that the smallest Fisher
  score on all instances is $4.77\times10^{-5}$, so the safeguard $\max\{\rho^{\top}w,10^{-15}\}$
  of the implementation is never active and the implemented $F_3$ is exactly~(41).
\item \textbf{Declarations.} The declaration of generative-AI use and the data-availability
  statement were updated, and a reproducibility package now accompanies the paper.
\end{enumerate}

\bigskip
\noindent We believe these revisions fully address the Reviewers' and Editor's comments
and have improved the clarity, rigor, and honesty of the manuscript. We thank the
Reviewers and the Editor again for their time.

\end{document}
